\documentclass[../main]{subfiles}
\begin{document}
\chapter{Motor Brushless}
\img{images/motor_brushless.jpg}{Motor Brushless}
\section{Introducción al Motor Brushless}

\subsection{Teoría y funcionamiento de los Motores Brushless}

Los motores brushless, o sin escobillas, son un tipo de motor eléctrico que se ha vuelto muy popular en diversas aplicaciones debido a su eficiencia, durabilidad y bajo mantenimiento. A diferencia de los motores con escobillas, los brushless no tienen contacto físico entre el rotor y el estator para la conmutación, lo que reduce significativamente el desgaste y aumenta la vida útil del motor.


En un motor brushless, el rotor es un imán permanente y el estator está compuesto por bobinas que se activan en secuencia mediante un controlador electrónico. Este controlador envía corrientes a las bobinas de manera que crean un campo magnético giratorio que interactúa con el campo magnético del rotor, haciendo que este gire.
\info{Principio de funcionamiento}{
La relación básica entre el voltaje $V$, la resistencia $R$, la inductancia $L$, la corriente $I$ y la velocidad angular $\omega$ en un motor brushless puede describirse con las siguientes fórmulas:
\begin{equation}
V = I \cdot R + L \cdot \frac{dI}{dt} + k_e \cdot \omega
\end{equation}
\begin{equation}
T = k_t \cdot I
\end{equation}
\begin{equation}
P = T \cdot \omega
\end{equation}

Donde:
\begin{itemize}
    \item $V$: Voltaje aplicado [V]
    \item $I$: Corriente eléctrica [A]
    \item $R$: Resistencia del bobinado [\Omega]
    \item $L$: Inductancia del bobinado [H]
    \item $\frac{dI}{dt}$: Derivada de la corriente con respecto al tiempo [A/s]
    \item $k_e$: Constante de fuerza contra-electromotriz [V/rad/s]
    \item $\omega$: Velocidad angular [rad/s]
    \item $T$: Par motor [Nm]
    \item $k_t$: Constante de par motor [Nm/A]
    \item $P$: Potencia mecánica [W]
\end{itemize}
}

Los motores brushless son esenciales en aplicaciones que requieren alta precisión y control de velocidad, como drones, vehículos eléctricos y equipos industriales.

\actividad{Responder el siguiente cuestionario}
\begin{enumerate}
\item ¿Qué es un motor brushless y en qué se diferencia de un motor con escobillas?
\item ¿Cuál es la función principal del controlador en un motor brushless?
\item Explica cómo se produce el movimiento en un motor brushless.
\item ¿Qué ventajas tienen los motores brushless sobre los motores con escobillas?
\item ¿Qué es la constante de fuerza contra-electromotriz ($k_e$) y cómo se relaciona con el voltaje y la velocidad angular?
\item Describe la relación entre el par motor ($T$) y la corriente ($I$) en un motor brushless.
\item ¿Cómo afecta la inductancia ($L$) del bobinado al funcionamiento de un motor brushless?
\item ¿Qué aplicaciones prácticas se benefician del uso de motores brushless?
\item ¿Cómo se puede medir la eficiencia de un motor brushless?
\item ¿Qué tipo de mantenimiento requieren los motores brushless en comparación con los motores con escobillas?

\end{enumerate}

\actividad{Resolver los siguientes problemas}
\begin{enumerate}
\item Calcular la corriente $I$ que circula por un motor brushless con una resistencia de bobinado $R = 0.5\ \Omega$, si se aplica un voltaje $V = 12\ V$ y la velocidad angular $\omega = 1000\ \text{rad/s}$, asumiendo que $k_e = 0.01\ \text{V/rad/s}$ y $L$ es despreciable.
\item Determinar el par motor $T$ de un motor brushless con una constante de par $k_t = 0.05\ \text{Nm/A}$ y una corriente $I = 10\ A$.
\item Un motor brushless entrega una potencia $P = 200\ W$ a una velocidad angular $\omega = 2000\ \text{rad/s}$. ¿Cuál es el par motor $T$ generado?
\item Si un motor brushless tiene una inductancia $L = 0.1\ H$ y la corriente aumenta a una tasa de $5\ \text{A/s}$, ¿qué voltaje adicional se induce debido a la inductancia?
\item Calcular el voltaje $V$ necesario para mantener una corriente $I = 8\ A$ en un motor brushless con $R = 0.2\ \Omega$, $L$ despreciable, $k_e = 0.02\ \text{V/rad/s}$ y $\omega = 1500\ \text{rad/s}$.
\item Un motor brushless tiene una resistencia $R = 1\ \Omega$ y se le aplica un voltaje de $24\ V$. Si la velocidad angular es $500\ \text{rad/s}$ y $k_e = 0.02\ \text{V/rad/s}$, ¿cuál es la corriente que circula por el motor?
\item Si un motor brushless con $k_t = 0.1\ \text{Nm/A}$ debe generar un par de $5\ Nm$, ¿qué corriente es necesaria?
\item Determinar la potencia mecánica $P$ de un motor brushless con par $T = 3\ Nm$ y velocidad angular $\omega = 2500\ \text{rad/s}$.
\item Calcular la velocidad angular $\omega$ de un motor brushless si se aplica un voltaje $V = 18\ V$, la corriente es $I = 6\ A$, $R = 0.3\ \Omega$, y $k_e = 0.015\ \text{V/rad/s}$.
\item Un motor brushless con una inductancia $L = 0.05\ H$ tiene una corriente que varía a una tasa de $10\ \text{A/s}$. ¿Qué voltaje se induce debido a esta variación de corriente?
\end{enumerate}

\end{document}
